% Generated semantic source for AI Almanach v3.
% Edit v3 directly after initial generation; do not overwrite
% editorial changes by rerunning the converter casually.

% ==== AI in E-Commerce, Marketing and Demand Forecasting =============================================================================
\chapter{AI in E-Commerce, Marketing and Demand Forecasting}

\chapterlead{How can personalisation create value without harmful feedback loops?}{Discover}{Recommend}{Experiment}{Govern}{commerce-marketing}

% ==== Section 1: AI in E-Commerce ===============================================================
\label{ch:ecommerce-marketing}

\begin{learningobjectives}
After working through this chapter you should be able to:
\begin{itemize}
  \item compute the sample size an A/B test needs to detect a given effect;
  \item explain why stopping a test when it reaches significance inflates the
        false-positive rate;
  \item distinguish a recommender that creates value from one that takes credit
        for it;
  \item name the feedback loop a recommender induces and its consequence for
        offline evaluation;
  \item state why offline recommender metrics disagree with online results.
\end{itemize}
\end{learningobjectives}

\begin{plainlanguage}
Commerce is the domain with the most data and the most self-deception. Every
click is recorded, every test is cheap, and almost every reported win is larger
than the real one --- because the metric was chosen after the result, the test
was stopped when it looked good, or the model was credited with a purchase that
would have happened anyway.
\end{plainlanguage}

\begin{engineernote}
The recurring error here is attributing to the model what belongs to the
counterfactual. A recommender shown on the homepage will always appear to
``drive'' revenue, because it is shown to people who came to buy. The only
honest measure is incrementality: what happened to a comparable group who did
not see it. Design the holdout before the system, because you cannot
reconstruct it afterwards.
\end{engineernote}

\section{AI in E-Commerce}

% ---- 16.1.1 Application Areas -----------------------------------------------------------------------
\subsection{Application Areas}


\begin{v3topic}{E-Commerce AI Landscape}
    E-commerce platforms apply AI across the full customer lifecycle\textsuperscript{\cite[][pp.~58--65]{laudonECommerce2020}}:
    \begin{itemize}
        \item \textbf{Retail}: product recommendations, dynamic pricing, inventory replenishment, visual search
        \item \textbf{Entertainment}: content personalisation, playlist generation, next-episode prediction
        \item \textbf{Advertising}: programmatic bidding, audience targeting, creative optimisation
        \item \textbf{Logistics}: last-mile routing, demand forecasting, warehouse automation
    \end{itemize}
    AI creates a \emph{data flywheel}: more users generate more behavioural data, enabling better models, attracting more users.
    \factvisual{
        \textbf{The flywheel has a governance condition.}
        More interactions can improve ranking and recommendations only when
        measurements represent durable user value.
        Optimising clicks alone can amplify popularity bias, manipulation, or
        short-term engagement at the expense of satisfaction.
    }{
        \begin{tikzpicture}[node distance=10mm]
            \node[almanach card] (users) {more users};
            \node[almanach card,right=of users] (events) {more events};
            \node[almanach accent card,below=of events] (models) {better models};
            \node[almanach card,left=of models] (value) {better experience};
            \draw[almanach arrow] (users) -- (events);
            \draw[almanach arrow] (events) -- (models);
            \draw[almanach arrow] (models) -- (value);
            \draw[almanach arrow] (value) -- (users);
            \node[font=\tiny\sffamily,text=AlmanachOchre,
                  align=center] at ($(users)!0.5!(models)$)
                  {measure\\long-term value};
        \end{tikzpicture}
        \visualcaption{A governed data flywheel. Feedback improves the system
          only when the objective measures the value the product intends to
          create.}{fig:ecommerce-data-flywheel}
    }

\end{v3topic}
\begin{v3topic}{Programmatic Advertising}
    \textbf{Real-time bidding (RTB)} auctions a single ad impression within ${\sim}100\,\mathrm{ms}$\textsuperscript{\cite[][pp.~445--452]{laudonECommerce2020}}.
    A \textbf{demand-side platform (DSP)} bids on behalf of advertisers; a \textbf{supply-side platform (SSP)} represents publishers.
    Models predict \textbf{click-through rate (CTR)} and \textbf{conversion rate (CVR)} from user context, page content, and historical behaviour.
\visualseparator
    \begin{tblr}{width=\linewidth,colspec={|X|X|}}
        \hline
        \textbf{Metric} & \textbf{Formula} \\
        \hline
        CTR & $\text{Clicks}/\text{Impressions}$ \\
        \hline
        CPM & Cost per 1000 impressions \\
        \hline
        ROAS & Revenue / Ad Spend \\
        \hline
    \end{tblr}
\captionof{table}{Programmatic advertising metrics and their formulas. Each is a ratio, so each can be improved by shrinking its denominator rather than by creating value.}
\label{tab:ecom-ad-metrics}

\end{v3topic}
\begin{v3topic}{IoT in Retail}
    \textbf{Smart retail} integrates sensor networks with AI analytics\textsuperscript{\cite[][pp.~70--75]{chaffeyDigitalBusinessECommerce2019}}:
    \begin{itemize}
        \item \textbf{RFID and shelf sensors}: real-time inventory visibility; trigger automated replenishment
        \item \textbf{Computer vision cameras}: dwell-time analysis, planogram compliance, cashierless checkout (Amazon Go)
        \item \textbf{Beacon technology}: proximity-based offers sent to shoppers' smartphones
        \item \textbf{Smart fitting rooms}: product information displayed via RFID tag reads
    \end{itemize}

\end{v3topic}
\begin{v3topic}{The Long Tail}
    In digital marketplaces, a vast number of niche products (the ``long tail'') collectively outperform a few popular items\textsuperscript{\cite[][pp.~3--8]{leskovecMiningMassiveDatasets2014}}.
    Traditional physical retail can stock only the most popular items; online platforms can offer millions.
    Recommender systems are the mechanism that makes the long tail economically viable: they connect users with relevant niche items that they would never discover by browsing.

\end{v3topic}

% ---- 16.1.2 Virtual Assistants and Search -----------------------------------------------------------------------
\subsection{Virtual Assistants and Search}


\begin{v3topic}{NLP-Powered Chatbots}
    E-commerce chatbots handle product queries, order tracking, and returns at scale\textsuperscript{\cite[][pp.~150--155]{chaffeyDigitalBusinessECommerce2019}}.
    Architecture: an intent classification model (fine-tuned transformer or BERT-based) maps user utterances to actions; a dialogue manager maintains context; a fulfilment layer queries backend systems.
    \textbf{Key metrics}: task completion rate, escalation rate, CSAT (customer satisfaction score).

\end{v3topic}
\begin{v3topic}{Voice Search}
    Voice queries (Alexa, Google, Siri) differ from text: they are longer, conversational, and often question-form\textsuperscript{\cite[][pp.~155--158]{chaffeyDigitalBusinessECommerce2019}}.
    Optimising for voice requires \textbf{featured snippets} (concise, direct answers) and structured data markup.
    Voice-first commerce: single-result responses mean the \emph{top-1 recommendation matters critically} --- ranking errors are immediately visible to the user.

\end{v3topic}
\begin{v3topic}{Visual Product Search}
    Visual search allows shoppers to photograph a product and retrieve visually similar items\textsuperscript{\cite[][pp.~160--163]{chaffeyDigitalBusinessECommerce2019}}.
    A CNN (e.g., ResNet or EfficientNet) encodes the query image into an embedding; approximate nearest-neighbour (ANN) search retrieves the closest product embeddings from the catalogue.
\visualseparator
    \textbf{Pipeline}: Query image $\xrightarrow{\text{CNN encoder}}$ $\mathbf{q} \in \mathbb{R}^d$ $\xrightarrow{\text{ANN index}}$ Top-$k$ product images.
    \textbf{Applications}: Pinterest Lens, Google Lens, ASOS visual search. Reduces the vocabulary gap: users can search for items they cannot name.

\end{v3topic}

% ---- 16.1.3 Dynamic Pricing -----------------------------------------------------------------------
\subsection{Dynamic Pricing}


\begin{v3topic}{Pricing Theory Foundations}
    The \textbf{price elasticity of demand} measures the percentage change in quantity demanded per percentage change in price\textsuperscript{\cite[][pp.~39--42]{laudonECommerce2020}}:
    \begin{equation}
        \varepsilon = \frac{\partial Q / Q}{\partial P / P} = \frac{\partial \ln Q}{\partial \ln P}
    \end{equation}
    $|\varepsilon| > 1$: elastic (price-sensitive); $|\varepsilon| < 1$: inelastic.
    \textbf{Willingness-to-pay (WTP)}: the maximum price a consumer accepts; price discrimination extracts consumer surplus by offering different prices to different segments.

\end{v3topic}
\begin{v3topic}{Bayesian Optimal Pricing}
    When true WTP is unknown, Bayesian inference maintains a posterior distribution over demand parameters\textsuperscript{\cite[][pp.~45--48]{laudonECommerce2020}}.
    Given a prior $p(\theta)$ and observed sales at historic prices, the posterior $p(\theta | \text{data})$ is updated via Bayes' theorem.
    The \textbf{optimal price} maximises expected revenue:
    \begin{equation}
        p^* = \arg\max_p \; p \cdot \mathbb{E}_\theta[Q(p, \theta)]
    \end{equation}
    Thompson Sampling balances \emph{exploration} (trying new prices) and \emph{exploitation} (using the current best estimate).

\end{v3topic}
\begin{v3topic}{RL-Based Dynamic Pricing}
    Pricing can be framed as an MDP: state $= $ inventory level, time-to-deadline, competitor prices; action $= $ posted price; reward $= $ revenue\textsuperscript{\cite[][pp.~55--60]{laudonECommerce2020}}.
    \textbf{Q-learning} and \textbf{policy gradient} methods learn pricing policies without assuming a demand model.
    \textbf{Surge pricing} (Uber, airlines) adjusts prices in real time to balance supply and demand.
    Key challenge: \emph{price fairness} --- identical items offered at different prices to different users raise ethical and legal concerns.

\end{v3topic}
\begin{v3topic}{Demand Elasticity Measurement}
    Estimating $\varepsilon$ from observational data is confounded by omitted variables (promotions, holidays)\textsuperscript{\cite[][p.~4]{wickDemandForecastingIndividual2021}}.
    \textbf{Approaches}:
    \begin{itemize}
        \item \textbf{Randomised price experiments}: A/B test different prices on matched user segments
        \item \textbf{Instrumental variables}: Use cost shocks (shipping, tariff changes) as instruments
        \item \textbf{Double ML} (Chernozhukov et al.): estimate nuisance functions with ML, then recover a causal price coefficient
    \end{itemize}

\end{v3topic}

% ---- 16.1.4 Ethics and Regulation -----------------------------------------------------------------------
\subsection{Ethics and Regulation}


\begin{v3topic}{GDPR in E-Commerce}
    The EU General Data Protection Regulation applies to any organisation processing EU residents' data\textsuperscript{\cite[][pp.~25--35]{itgovernanceprivacyteamEUGDPR2020}}:
    \begin{itemize}
        \item \textbf{Lawful basis}: consent or legitimate interest for personalisation and advertising cookies
        \item \textbf{Data minimisation}: collect only what is strictly necessary; avoid re-identification of anonymised data
        \item \textbf{Right to explanation}: profiling and automated decisions must be explainable to data subjects
        \item \textbf{Data portability}: users can request their data in machine-readable format
    \end{itemize}

\end{v3topic}
\begin{v3topic}{Algorithmic Fairness and Dark Patterns}
    \textbf{Algorithmic pricing fairness}: dynamic pricing must not discriminate based on protected characteristics (race, gender) embedded in correlated features such as neighbourhood\textsuperscript{\cite[][pp.~80--85]{laudonECommerce2020}}.
    \textbf{Dark patterns} are UI/UX designs that exploit cognitive biases to nudge users into unintended actions (e.g., hidden unsubscribe flows, pre-ticked consent boxes, countdown timers).
    The EU Digital Services Act (DSA, 2022) explicitly prohibits dark patterns targeting consumers.

\end{v3topic}

% ==== Section 2: Marketing Analytics ===============================================================
\section{Marketing Analytics}

% ---- 16.2.1 Foundations -----------------------------------------------------------------------
\subsection{Foundations}


\begin{v3topic}{Marketing Mix (4Ps)}
    The classical \textbf{marketing mix} describes four controllable decision variables\textsuperscript{\cite[][pp.~18--24]{chaffeyDigitalBusinessECommerce2019}}:
    \begin{itemize}
        \item \textbf{Product}: features, quality, branding, packaging
        \item \textbf{Price}: pricing strategy, discounts, payment terms
        \item \textbf{Place}: distribution channels (own site, marketplace, omnichannel)
        \item \textbf{Promotion}: advertising, content marketing, email, SEO/SEM
    \end{itemize}
    Digital channels have added \textbf{4Cs} (Customer, Cost, Convenience, Communication) as a consumer-centric counterpart.

\end{v3topic}
\begin{v3topic}{Customer Journey and Attribution}
    The \textbf{customer journey} maps touchpoints from awareness to purchase and retention\textsuperscript{\cite[][pp.~92--98]{chaffeyDigitalBusinessECommerce2019}}.
    \textbf{Attribution models} assign credit for a conversion across touchpoints:
    \begin{itemize}
        \item \textbf{Last-click}: 100\% credit to final touchpoint (over-values retargeting)
        \item \textbf{First-click}: 100\% credit to first touchpoint (over-values awareness)
        \item \textbf{Linear}: equal credit to all touchpoints
        \item \textbf{Data-driven (algorithmic)}: Shapley values from game theory allocate credit proportionally to each channel's marginal contribution
    \end{itemize}

\end{v3topic}

% ---- 16.2.2 Descriptive Analytics -----------------------------------------------------------------------
\subsection{Descriptive Analytics}


\begin{v3topic}{RFM Customer Segmentation}
    \textbf{RFM analysis} scores customers on three dimensions to identify high-value segments\textsuperscript{\cite[][pp.~110--115]{chaffeyDigitalBusinessECommerce2019}}:
    \begin{itemize}
        \item \textbf{Recency (R)}: days since last purchase --- recent buyers are more likely to respond
        \item \textbf{Frequency (F)}: number of purchases in a period --- loyal customers are more profitable
        \item \textbf{Monetary (M)}: total spend --- high-value customers deserve premium service
    \end{itemize}
    Customers are ranked 1--5 on each dimension; a score of 555 identifies champions. RFM segments guide differentiated CRM campaigns.

\end{v3topic}
\begin{v3topic}{Market Basket Analysis}
    \textbf{Association rule mining} discovers co-purchase patterns in transaction data\textsuperscript{\cite[][pp.~217--221]{leskovecMiningMassiveDatasets2014}}.
    For a rule $A \Rightarrow B$:
    \begin{equation}
        \text{support}(A \cup B) = \frac{|\{t : A \cup B \subseteq t\}|}{|T|}
    \end{equation}
    \begin{equation}
        \text{confidence}(A \Rightarrow B) = \frac{\text{support}(A \cup B)}{\text{support}(A)}
    \end{equation}
    \begin{equation}
        \text{lift}(A \Rightarrow B) = \frac{\text{confidence}(A \Rightarrow B)}{\text{support}(B)}
    \end{equation}
    Lift $> 1$ indicates a genuine association beyond chance.

\end{v3topic}
\begin{v3topic}{Apriori Algorithm}
    The \textbf{Apriori algorithm} efficiently mines frequent itemsets by exploiting the \emph{anti-monotone property}: any superset of an infrequent itemset is also infrequent\textsuperscript{\cite[][pp.~222--226]{leskovecMiningMassiveDatasets2014}}.
\visualseparator
    \textbf{Procedure}:
    \begin{enumerate}
        \item Find all frequent 1-itemsets (items with support $\geq$ \textit{min\_sup})
        \item Generate candidate $k$-itemsets by joining frequent $(k-1)$-itemsets
        \item Prune candidates with any infrequent $(k-1)$-subset
        \item Repeat until no new frequent itemsets found
    \end{enumerate}
    \textbf{FP-Growth} improves efficiency by compressing transactions into an FP-tree, avoiding repeated database scans.

\end{v3topic}
\begin{v3topic}{SEO and Search Analytics}
    \textbf{Search engine optimisation (SEO)} improves organic visibility\textsuperscript{\cite[][pp.~130--136]{chaffeyDigitalBusinessECommerce2019}}.
    Key signals for search ranking: relevance (keyword match, semantic similarity via BERT-based models), authority (PageRank, backlinks), user experience (Core Web Vitals, bounce rate, dwell time).
    \textbf{Search analytics}: query performance reports reveal which terms drive traffic, click-through rates per keyword, and content gap opportunities.
    \textbf{SEM}: paid search (Google Ads) complements organic reach via automated bidding strategies optimised for conversions.

\end{v3topic}

% ---- 16.2.3 Predictive Analytics -----------------------------------------------------------------------
\subsection{Predictive Analytics}


\begin{v3topic}{Customer Churn Prediction}
    \textbf{Churn} is the rate at which customers stop doing business with a company\textsuperscript{\cite[][pp.~145--150]{chaffeyDigitalBusinessECommerce2019}}.
    Prediction pipeline:
    \begin{enumerate}
        \item \textbf{Define churn}: no purchase for $N$ days (contractual) or declining engagement (non-contractual)
        \item \textbf{Feature engineering}: RFM scores, session frequency, support tickets, time-since-last-login
        \item \textbf{Model}: gradient boosted trees (XGBoost/LightGBM) or survival models (Cox proportional hazards)
        \item \textbf{Action}: target high-propensity churners with retention offers
    \end{enumerate}

\end{v3topic}
\begin{v3topic}{Customer Lifetime Value (CLV)}
    \textbf{CLV} estimates the total net profit attributable to a customer over their relationship with the firm\textsuperscript{\cite[][pp.~152--157]{chaffeyDigitalBusinessECommerce2019}}.
    The basic discounted cash flow formulation:
    \begin{equation}
        \text{CLV} = \sum_{t=1}^{T} \frac{m_t \cdot r_t}{(1+d)^t} - \text{CAC}
    \end{equation}
    where $m_t$ = margin in period $t$, $r_t$ = retention probability, $d$ = discount rate, CAC = customer acquisition cost.
    The \textbf{Pareto/NBD model} (Schmittlein et al.) estimates CLV for non-contractual settings using purchase history.

\end{v3topic}
\begin{v3topic}{Sales Forecasting}
    Accurate sales forecasts drive inventory planning, staffing, and budgeting\textsuperscript{\cite[][pp.~4--8]{hyndmanForecastingPrinciples2021}}.
    \textbf{Methods hierarchy}:
    \begin{itemize}
        \item \textbf{Statistical}: ARIMA, Exponential Smoothing (ETS), seasonal decomposition
        \item \textbf{ML}: gradient boosted trees with lag features, calendar effects, price features
        \item \textbf{Deep learning}: DeepAR (probabilistic LSTM), N-BEATS, temporal fusion transformer
        \item \textbf{Ensemble}: weighted average of multiple models; meta-learner approach
    \end{itemize}
    Probabilistic forecasts (prediction intervals) are preferable to point estimates for inventory decisions.

\end{v3topic}
\begin{v3topic}{Propensity Models}
    A \textbf{propensity model} estimates the probability that a customer will perform a target action (purchase, upgrade, open email)\textsuperscript{\cite[][pp.~162--167]{chaffeyDigitalBusinessECommerce2019}}.
    Training: binary classification on historical data; features include browsing behaviour, demographics, product interest signals.
    \textbf{Uplift modelling} (causal ML) goes further: it estimates the \emph{incremental} effect of a treatment (e.g., sending a coupon) on conversion probability, avoiding wasted spend on customers who would convert anyway.

\end{v3topic}

% ---- 16.2.4 Prescriptive Analytics -----------------------------------------------------------------------
\subsection{Prescriptive Analytics}


\begin{v3topic}{A/B Testing and Experiments}
    \textbf{Randomised controlled experiments} are the gold standard for evaluating marketing interventions\textsuperscript{\cite[][pp.~178--183]{chaffeyDigitalBusinessECommerce2019}}.
    Procedure: randomly split users into control (A) and treatment (B); measure the metric of interest; apply a two-sample hypothesis test.
    \textbf{Minimum detectable effect (MDE)}: determines required sample size $n \propto (\sigma / \delta)^2$, where $\sigma$ is variance and $\delta$ is the smallest meaningful effect.
    Pitfalls: peeking (stopping early), novelty effects, interaction effects between concurrent experiments.

\end{v3topic}
\begin{v3topic}{Multi-Armed Bandits for Campaigns}
    \textbf{Multi-armed bandit (MAB)} algorithms allocate traffic adaptively across campaign variants, balancing exploration and exploitation\textsuperscript{\cite[][pp.~458--462]{suttonReinforcementLearningIntroduction2018}}:
    \begin{itemize}
        \item \textbf{$\varepsilon$-greedy}: with probability $\varepsilon$ explore randomly; otherwise exploit the best-known variant
        \item \textbf{UCB}: select variant with highest upper confidence bound $\hat{\mu}_i + \sqrt{2\ln t / n_i}$
        \item \textbf{Thompson Sampling}: sample from Beta posterior; select highest sample
    \end{itemize}
    MABs converge faster than fixed A/B tests and minimise regret during the experiment itself.

\end{v3topic}
\begin{v3topic}{Upselling, Cross-selling and Targeting}
    \textbf{Upselling}: recommending a higher-tier product to a customer already considering a purchase.
    \textbf{Cross-selling}: recommending complementary products (``Customers also bought\ldots'').
    Both are powered by association rules and collaborative filtering models\textsuperscript{\cite[][pp.~185--190]{chaffeyDigitalBusinessECommerce2019}}.
    \textbf{Audience targeting}: lookalike models identify new prospects who resemble existing high-value customers; features are matched via embedding similarity in ad platform user spaces.

\end{v3topic}
\begin{v3topic}{Closed-Loop vs.\ Human-in-the-Loop}
    \textbf{Closed-loop marketing automation}: ML models select, personalise, and send communications without human review; feedback (opens, clicks, conversions) retrains models automatically\textsuperscript{\cite[][pp.~200--205]{chaffeyDigitalBusinessECommerce2019}}.
    \textbf{Human-in-the-loop}: marketers review ML recommendations before deployment; essential for brand safety, legal compliance, and novel contexts where models lack training data.
    \textbf{Omnichannel}: unifying customer data and messaging across web, mobile, email, in-store, and social for a consistent journey.

\end{v3topic}

% ==== Section 3: Recommender Systems ===============================================================
\section{Recommender Systems}

% ---- 16.3.1 Foundations -----------------------------------------------------------------------
\subsection{Foundations}


\begin{v3topic}{History and Impact}
    Recommender systems emerged in the mid-1990s with GroupLens (Usenet news) and Amazon's item-based filtering\textsuperscript{\cite[][pp.~1--5]{aggarwalRecommenderSystems2016}}.
    The \textbf{Netflix Prize} (2006--2009) catalysed algorithmic research: a \$1M prize for improving RMSE by 10\% over Cinematch generated breakthrough matrix factorisation techniques.
    Today, recommenders drive $\sim$35\% of Amazon's revenue and $\sim$80\% of Netflix viewing time.

\end{v3topic}
\begin{v3topic}{User--Item Matrix and Feedback Types}
    The central data structure is the \textbf{user--item interaction matrix} $\mathbf{R} \in \mathbb{R}^{m \times n}$, where $m$ = users, $n$ = items\textsuperscript{\cite[][pp.~8--13]{aggarwalRecommenderSystems2016}}.
\visualseparator
    \begin{tblr}{width=\linewidth,colspec={|X|X|}}
        \hline
        \textbf{Explicit feedback} & \textbf{Implicit feedback} \\
        \hline
        Star ratings (1--5) & Clicks, views, dwell time \\
        \hline
        Thumbs up/down & Add to cart, purchase \\
        \hline
        Sparse, noisy & Dense, but ambiguous \\
        \hline
        User-intentional & Behavioural proxy \\
        \hline
    \end{tblr}
\captionof{table}{Explicit against implicit feedback. Implicit feedback is abundant and ambiguous: a click is not an endorsement, and an absent click is not a rejection.}
\label{tab:ecom-feedback-types}
    Most production systems rely on implicit feedback because it is abundant, though it does not distinguish preference from mere exposure.

\end{v3topic}
\begin{v3topic}{Levels of Personalisation}
    Recommender systems exist on a personalisation spectrum\textsuperscript{\cite[][pp.~15--18]{jannachRecommenderSystemsIntroduction2010}}:
    \begin{itemize}
        \item \textbf{Non-personalised}: editorial ``bestsellers'' or ``trending'' lists --- no user model required
        \item \textbf{Segment-based}: same recommendations for a user group (age, geography) --- coarse but scalable
        \item \textbf{Collaborative}: personalised based on similar users' behaviour --- requires interaction data
        \item \textbf{Content-based}: personalised based on item attributes and user profile --- effective for new items
        \item \textbf{Hybrid}: combines multiple signals --- best accuracy in practice
    \end{itemize}

\end{v3topic}
\begin{v3topic}{Evaluation Metrics}
    Offline evaluation metrics for recommenders\textsuperscript{\cite[][pp.~20--28]{aggarwalRecommenderSystems2016}}:
    \begin{itemize}
        \item \textbf{RMSE/MAE}: for explicit-rating prediction tasks
        \item \textbf{Precision@$k$}: fraction of recommended items that are relevant
        \item \textbf{Recall@$k$}: fraction of relevant items that appear in top-$k$
        \item \textbf{NDCG@$k$}: normalised discounted cumulative gain, rewards rank-ordering of relevant items
        \item \textbf{Diversity}: average pairwise dissimilarity in top-$k$ list
        \item \textbf{Coverage}: fraction of the catalogue that the system recommends across users
    \end{itemize}
    Offline metrics are necessary but not sufficient: \textbf{online A/B tests} measure business impact (CTR, conversion, revenue).

\end{v3topic}

% ---- 16.3.2 Collaborative Filtering -----------------------------------------------------------------------
\subsection{Collaborative Filtering}


\begin{v3topic}{User-Based Collaborative Filtering}
    \textbf{User-based CF} predicts user $u$'s rating for item $i$ by aggregating ratings from similar users\textsuperscript{\cite[][pp.~31--36]{aggarwalRecommenderSystems2016}}:
    \begin{equation}
        \hat{r}_{ui} = \bar{r}_u + \frac{\sum_{v \in N(u)} \text{sim}(u,v)\,(r_{vi} - \bar{r}_v)}{\sum_{v \in N(u)} |\text{sim}(u,v)|}
    \end{equation}
    \textbf{Cosine similarity} between user rating vectors:
    \begin{equation}
        \text{sim}(u,v) = \frac{\mathbf{r}_u \cdot \mathbf{r}_v}{\|\mathbf{r}_u\|\,\|\mathbf{r}_v\|}
    \end{equation}
    \textbf{Pearson correlation} accounts for rating scale differences between users.

\end{v3topic}
\begin{v3topic}{Item-Based Collaborative Filtering}
    \textbf{Item-based CF} (Amazon's original algorithm) computes item--item similarities from user rating patterns\textsuperscript{\cite[][pp.~38--42]{aggarwalRecommenderSystems2016}}:
    \begin{equation}
        \hat{r}_{ui} = \frac{\sum_{j \in N(i)} \text{sim}(i,j)\,r_{uj}}{\sum_{j \in N(i)} |\text{sim}(i,j)|}
    \end{equation}
    Item similarities are \textbf{pre-computed offline} and stable over time; only the personalised combination at query time is online.
    \textbf{Advantage over user-based}: items change less frequently than users; scales better for large user bases.

\end{v3topic}
\begin{v3topic}{Matrix Factorisation}
    \textbf{Matrix factorisation (MF)} decomposes the interaction matrix $\mathbf{R} \approx \mathbf{P}\mathbf{Q}^\top$, where $\mathbf{P} \in \mathbb{R}^{m \times k}$ contains \emph{user latent factors} and $\mathbf{Q} \in \mathbb{R}^{n \times k}$ contains \emph{item latent factors}\textsuperscript{\cite[][pp.~30--37]{korenMatrixFactorizationTechniques2009}}.
    The predicted rating:
    \begin{equation}
        \hat{r}_{ui} = \mathbf{p}_u^\top \mathbf{q}_i + b_u + b_i + \mu
    \end{equation}
    where $b_u$, $b_i$ are user/item biases and $\mu$ is the global mean.
    \textbf{SVD} learns factors by minimising regularised squared error via stochastic gradient descent:
    \begin{equation}
        \min_{\mathbf{P},\mathbf{Q}} \sum_{(u,i)\in\mathcal{K}} \bigl(r_{ui} - \hat{r}_{ui}\bigr)^2 + \lambda\bigl(\|\mathbf{p}_u\|^2 + \|\mathbf{q}_i\|^2\bigr)
    \end{equation}

\end{v3topic}
\begin{v3topic}{ALS for Implicit Feedback}
    \textbf{Alternating Least Squares (ALS)} is the standard approach for implicit feedback\textsuperscript{\cite[][pp.~40--45]{korenMatrixFactorizationTechniques2009}}.
    A binary confidence matrix $c_{ui} = 1 + \alpha \cdot \text{count}_{ui}$ encodes interaction strength.
    ALS fixes $\mathbf{Q}$ and solves for each $\mathbf{p}_u$ in closed form (least-squares), then swaps.
    Because the closed-form solution is the same for all users given $\mathbf{Q}$, it parallelises efficiently for large-scale systems.

\end{v3topic}
\begin{v3topic}{Bayesian Personalised Ranking (BPR)}
    \textbf{BPR} optimises for \emph{pairwise ranking} rather than rating prediction\textsuperscript{\cite[][p.~1]{rendleBPRBayesianPersonalized2009}}: the model should score observed (positive) interactions higher than unobserved ones.
    For each triple $(u, i^+, j^-)$ (user $u$, observed item $i^+$, unobserved item $j^-$):
    \begin{equation}
        \mathcal{L}_{\text{BPR}} = -\sum_{(u,i^+,j^-)} \ln\sigma\!\bigl(\hat{r}_{ui^+} - \hat{r}_{uj^-}\bigr) + \lambda \|\Theta\|^2
    \end{equation}
    BPR is model-agnostic: it can be applied on top of MF, neural networks, or any scoring function.
    \textbf{Key insight}: for e-commerce, the task is \emph{ranking} (what to show at position 1), not \emph{rating prediction} (what score to assign).

\end{v3topic}

% ---- 16.3.3 Content-Based Filtering -----------------------------------------------------------------------
\subsection{Content-Based Filtering}


\begin{v3topic}{Feature Extraction for Items}
    Content-based filtering builds an item profile from attributes and recommends items similar to those the user has liked\textsuperscript{\cite[][pp.~65--70]{aggarwalRecommenderSystems2016}}.
    \textbf{Text}: TF-IDF or sentence embeddings (BERT) of product descriptions and reviews.
    \textbf{Images}: CNN feature vectors (ResNet last layer) for visual products (fashion, home decor).
    \textbf{Structured}: categorical encodings (genre, brand, category), numerical features (price, rating).
    A \textbf{user profile} is built as a weighted average of item features rated positively; recommendations are items with highest cosine similarity to the profile.

\end{v3topic}
\begin{v3topic}{Factorisation Machines}
    \textbf{Factorisation Machines (FM)} generalise matrix factorisation to arbitrary feature vectors\textsuperscript{\cite[][p.~995]{rendleFactorizationMachines2010}}:
    \begin{equation}
        \hat{y}(\mathbf{x}) = w_0 + \sum_i w_i x_i + \sum_{i<j} \langle \mathbf{v}_i, \mathbf{v}_j \rangle x_i x_j
    \end{equation}
    The pairwise interaction term $\langle \mathbf{v}_i, \mathbf{v}_j \rangle$ captures feature interactions (e.g., user $\times$ item, item $\times$ category) via low-rank factorisation of the interaction matrix.
    FM unifies MF and content-based approaches: user, item, and context features can all be included in $\mathbf{x}$.

\end{v3topic}

% ---- 16.3.4 Hybrid Systems -----------------------------------------------------------------------
\subsection{Hybrid Recommender Systems}


\begin{v3topic}{Hybrid Approaches}
    Hybrid systems combine collaborative and content-based methods to mitigate their individual weaknesses\textsuperscript{\cite[][pp.~80--87]{jannachRecommenderSystemsIntroduction2010}}:
    \begin{itemize}
        \item \textbf{Weighted hybrid}: linear combination of scores from multiple recommenders
        \item \textbf{Switching}: use content-based for new users (cold-start), switch to CF once enough interactions exist
        \item \textbf{Cascade}: use one method to filter candidates, another to re-rank
        \item \textbf{Feature augmentation}: output of one model used as input feature to another
        \item \textbf{Meta-level}: user model of one system used as input to another
    \end{itemize}

\end{v3topic}
\begin{v3topic}{The Cold-Start Problem}
    \textbf{Cold start} occurs when a new user or new item has no interaction history\textsuperscript{\cite[][pp.~88--92]{aggarwalRecommenderSystems2016}}:
    \begin{itemize}
        \item \textbf{New user}: ask for explicit preferences (onboarding quiz); use demographic defaults; show popularity-based recommendations
        \item \textbf{New item}: use content features (no interaction needed); boost new items with controlled exploration budget
        \item \textbf{Contextual}: use session context (query, referrer, device) as surrogate for user model
    \end{itemize}
    Cold start is why purely collaborative systems fail in practice without content-based fallbacks.

\end{v3topic}

% ---- 16.3.5 Large-Scale and Modern Approaches -----------------------------------------------------------------------
\subsection{Large-Scale and Modern Approaches}


\begin{v3topic}{Two-Stage Retrieval Architecture}
    Large-scale recommenders (billions of users, millions of items) separate the problem into two stages\textsuperscript{\cite[][pp.~95--100]{aggarwalRecommenderSystems2016}}:
    \begin{enumerate}
        \item \textbf{Retrieval (candidate generation)}: fast approximate nearest-neighbour search returns $\mathcal{O}(100)$ candidates from the full catalogue
        \item \textbf{Ranking}: a complex model (Wide\&Deep, transformer) scores the candidates and selects the final top-$k$
    \end{enumerate}
\visualseparator
    \textbf{Why two stages?} Expensive ranking models cannot evaluate millions of items per request; fast retrieval narrows the search space.
    \textbf{Two-tower model}: a query tower encodes the user/context; an item tower encodes item features; the inner product $\mathbf{q}_u^\top \mathbf{e}_i$ is the retrieval score.

\end{v3topic}
\begin{v3topic}{Approximate Nearest Neighbour (ANN)}
    \textbf{FAISS} (Facebook AI Similarity Search) provides GPU-accelerated exact and approximate nearest-neighbour search over embedding vectors\textsuperscript{\cite[][p.~535]{liFaissLibraryEfficient2019}}.
    Key FAISS indices:
    \begin{itemize}
        \item \textbf{IVF} (Inverted File): partitions the space into Voronoi cells; search only a subset of cells
        \item \textbf{HNSW} (Hierarchical Navigable Small World): graph-based index with $O(\log n)$ query complexity
        \item \textbf{PQ} (Product Quantisation): compresses vectors to reduce memory; trades recall for speed
    \end{itemize}

\end{v3topic}
\begin{v3topic}{Wide \& Deep / Neural CF}
    \textbf{Wide \& Deep} (Google, 2016) combines a \emph{wide} linear model (for memorisation of frequent patterns) with a \emph{deep} MLP (for generalisation)\textsuperscript{\cite[][p.~7]{chengWideDeepLearning2016}}.
    \textbf{Neural Collaborative Filtering (NCF)} (He et al., 2017) replaces the inner product in MF with an MLP to learn arbitrary user--item interactions\textsuperscript{\cite[][p.~173]{heNeuralCollaborativeFiltering2017}}:
    \begin{equation}
        \hat{r}_{ui} = \sigma\!\bigl(\mathbf{h}^\top \text{MLP}\bigl([\mathbf{p}_u \| \mathbf{q}_i]\bigr)\bigr)
    \end{equation}
    NCF achieves state-of-the-art ranking performance on implicit feedback benchmarks.

\end{v3topic}
\begin{v3topic}{RL and Causal Approaches}
    \textbf{RL-based recommenders} model the recommendation session as an MDP: state = user history; action = item to recommend; reward = long-term engagement or revenue\textsuperscript{\cite[][pp.~470--475]{suttonReinforcementLearningIntroduction2018}}.
    Optimising for immediate clicks creates filter bubbles; RL can explicitly trade short-term engagement for diversity and long-term retention.
    \textbf{Causal recommenders} address \emph{exposure bias}: items that were never shown cannot appear in training data as negatives; inverse propensity scoring (IPS) re-weights observations to correct for the logging policy\textsuperscript{\cite[][pp.~22--30]{pearlCausalInferenceStatistics2016}}.

\end{v3topic}
\begin{v3topic}{Multi-Stakeholder Recommenders}
    Real platforms serve multiple principals simultaneously\textsuperscript{\cite[][pp.~105--110]{aggarwalRecommenderSystems2016}}:
    \begin{itemize}
        \item \textbf{Users}: want relevant, diverse, non-repetitive recommendations
        \item \textbf{Producers/sellers}: want fair visibility for their items (exposure fairness)
        \item \textbf{Platform}: maximises revenue, engagement, and regulatory compliance
    \end{itemize}
    \textbf{Multi-objective optimisation}: Pareto-optimal trade-offs between user satisfaction and producer fairness; constrained re-ranking adds minimum exposure constraints for small sellers without critically degrading user utility.

\end{v3topic}

% ==== Section 4: Further Reading ===============================================================
%
\section{Deep Dive: Feedback Loops and Multi-Sided Objectives}
\begin{v3topic}{Logged Data Reflects the Old Policy}
Clicks and purchases are observed only for items exposed by an earlier ranking
and interface.
Position, availability, and selection bias confound naive supervised targets.
Randomised exploration or carefully justified propensity methods help estimate
counterfactual performance.

\end{v3topic}
\begin{v3topic}{Delayed and Interfering Outcomes}
Campaigns can change later conversion, returns, fatigue, and word of mouth.
One user's exposure can affect inventory, prices, sellers, and other users, so
standard independent-unit A/B assumptions may fail.

\end{v3topic}
\begin{v3topic}{Balance the Marketplace}
Evaluate relevance, conversion, margin, diversity, novelty, long-term retention,
producer exposure, support burden, and safety constraints.
Define which objectives are hard constraints and which can trade off.
Monitor feedback loops that make the training data increasingly reflect the
model's own past choices.
\end{v3topic}


\begin{workedexample}[How long must the A/B test run?]
Baseline conversion $p_0 = 4\,\%$. You want to detect a relative improvement of
$5\,\%$, i.e.\ $p_1 = 4.2\,\%$, at $\alpha = 0.05$ and $80\,\%$ power.

The standard approximation for two proportions is
\[
  n \approx \frac{2\,\bigl(z_{\alpha/2}+z_{\beta}\bigr)^2\,\bar p\,(1-\bar p)}
                 {(p_1-p_0)^2},
  \qquad \bar p = 0.041 .
\]
With $z_{\alpha/2} = 1.96$ and $z_\beta = 0.84$, so
$(1.96+0.84)^2 = 7.85$:
\[
  n \approx \frac{2(7.85)(0.041)(0.959)}{(0.002)^2}
    = \frac{0.6174}{0.000004}
    \approx 154{,}000 \text{ per arm.}
\]

\emph{Roughly $308{,}000$ visitors in total.} At $10{,}000$ visitors a day the
test needs about a month.

\emph{The lesson in the denominator.} The required sample scales as
$1/(\text{effect size})^2$. Halving the effect you want to detect --- to a
$2.5\,\%$ relative lift --- quadruples the sample to about $616{,}000$ per arm.
Small effects are not slightly harder to measure; they are quadratically
harder. If you cannot reach the sample size, the honest conclusion is that the
experiment cannot answer the question, not that the result was inconclusive.
\end{workedexample}

\begin{cautionbox}[Peeking turns a 5\,\% error rate into something much worse]
Checking a running test and stopping as soon as $p < 0.05$ appears is the most
common error in commercial experimentation. Under a true null, the $p$-value
wanders; given enough looks it will cross $0.05$ eventually. Testing
continuously drives the false-positive rate far above the nominal $5\,\%$ ---
with frequent peeking it can exceed $30\,\%$.

The fixes: fix the sample size in advance and look once; or use a method
designed for continuous monitoring --- sequential testing with alpha spending, or
a Bayesian approach with a pre-declared decision rule. What you may not do is
use a fixed-horizon $p$-value and stop whenever you like. This, combined with
running many simultaneous tests without correction, is why so many reported
uplifts do not replicate.
\end{cautionbox}

\begin{cautionbox}[Recommenders grade their own homework]
A recommender trained on logged interactions learns from data it caused. Items
it surfaced get clicked and become training signal; items it never surfaced
generate no data and are learned to be uninteresting. The loop is
self-confirming, and offline metrics computed on logged data reward exactly
this behaviour --- which is why offline ranking scores routinely fail to predict
online results.

Two partial defences: reserve a small randomised traffic slice that shows
unranked or randomised items, giving unbiased data for evaluation; and measure
incrementality against a holdout that sees no recommendations at all. Both cost
short-term revenue, and both are the only way to know whether the system is
creating value or redistributing it.
\end{cautionbox}

\begin{practicallab}[Measure incrementality, not attribution]
\textbf{Goal.} Find out what a recommender is actually worth.

\textbf{Protocol.}
\begin{enumerate}
  \item Define the metric before anything else --- revenue per visitor, not
        click-through --- and pre-register it in writing.
  \item Compute the required sample size for the smallest effect worth acting
        on. If it is unreachable, stop and say so.
  \item Hold out a random slice of users who see no recommendations. This is
        the only group that answers the question.
  \item Run to the pre-computed sample size. Do not look at significance before
        then; monitor only for outages.
  \item Report the effect size with a confidence interval, alongside the
        $p$-value.
  \item Compute what naive attribution \emph{would} have claimed, and compare
        with the incremental result.
\end{enumerate}

\textbf{Deliverable.} The pre-registration, the required $n$, the incremental
lift with its interval, and the attribution figure beside it.
\textbf{The point:} step 6 typically shows attribution overstating the effect
by a factor of two or more.
\end{practicallab}

\begin{checkpoint}
A recommender is credited with $18\,\%$ of revenue. Give three reasons this
number overstates its value, and describe the single experiment that would
establish the real one.
\end{checkpoint}

\chapterreview{Discover}{Recommend}{Experiment}{Govern}

\section{Further Reading}

\begin{itemize}
    \item \fullcite{aggarwalRecommenderSystems2016} --- Comprehensive graduate-level textbook covering all recommender paradigms: neighbourhood methods, matrix factorisation, deep learning, and context-aware systems.
    \item \fullcite{jannachRecommenderSystemsIntroduction2010} --- Accessible introduction to recommender systems with practical examples; covers collaborative, content-based, and hybrid methods.
    \item \fullcite{ricciRecommenderSystemsHandbook2015} --- Edited handbook with 30+ chapters from leading researchers; the reference for advanced topics including cross-domain and group recommenders.
    \item \fullcite{leskovecMiningMassiveDatasets2014} --- Chapters 9 (recommendation systems) and 6 (frequent itemsets / Apriori) are directly relevant; freely available online.
    \item \fullcite{laudonECommerce2020} --- Broad e-commerce textbook; covers digital business models, advertising technology, pricing, and regulatory landscape.
    \item \fullcite{chaffeyDigitalBusinessECommerce2019} --- Practitioner-oriented digital business guide; strong on marketing analytics, customer journey, and omnichannel strategy.
    \item \fullcite{itgovernanceprivacyteamEUGDPR2020} --- Implementation guide for GDPR compliance; directly applicable to data collection, profiling, and personalisation in e-commerce.
    \item \fullcite{hyndmanForecastingPrinciples2021} --- The standard reference for time-series forecasting methods; freely available online; essential for demand and sales forecasting.
    \item \fullcite{korenMatrixFactorizationTechniques2009} --- Seminal paper on matrix factorisation for recommender systems; introduces SVD with biases; highly readable introduction to latent factor models.
    \item \fullcite{pearlCausalInferenceStatistics2016} --- Primer on causal inference; essential background for uplift modelling, debiased recommenders, and causal attribution in marketing.
\end{itemize}
